\documentclass[../../main.tex]{subfiles}

\graphicspath{{../../image/}{../../figure/}}

\begin{document}

\section{习题}\label{section:ch16-习题}

\begin{exercise}\label{exercise:ch16-01}
（$\star\star$）使用归纳法证明，将数据投影到 $M$ 维子空间以最大化投影数据的方差所定义的线性投影是由数据协方差矩阵 $\bm{S}$[\eqref{align:ch16-协方差}]的 $M$ 个特征向量定义的，对应于 $M$ 个最大特征值。在\cref{subsection:ch16-最大方差表述}中，这个结果已经在 $M=1$ 的情况下得到证明。假设该结果对某个一般的 $M$ 值成立，证明它因此对 $M+1$ 也成立。为此，首先将投影数据的方差关于定义新方向的向量 $\bm{u}_{M+1}$ 的导数设置为零。这应该在满足约束条件的情况下完成，即 $\bm{u}_{M+1}$ 与现有向量 $\bm{u}_1,\cdots,\bm{u}_M$ 正交，并且已归一化至单位长度。使用 Lagrange 乘子保证这些约束，然后利用向量 $\bm{u}_1,\cdots,\bm{u}_M$ 的正交性质证明新向量 $\bm{u}_{M+1}$ 是 $\bm{S}$ 的一个特征向量。最后证明如果选择与第 $M+1$ 个特征值相对应的特征向量（特征值按降序排列），则方差达到最大。
\end{exercise}

\begin{exercise}\label{exercise:ch16-02}
（$\star\star$）证明在满足正交性质[\eqref{align:ch16-正交基}]的情况下，PCA 误差度量 $J$[\eqref{align:ch16-误差函数}]的最小值是在 $\bm{u}_j$ 为数据协方差矩阵 $\bm{S}$ 的特征向量时获得的。为此，引入一个 Lagrange 乘子矩阵 $\bm{H}$，每个约束对应一个 Lagrange 乘子，因此修改后的误差度量 $\tilde{J}$ 用矩阵表示为
\begin{align}
  \tilde{J}=\operatorname{tr}\left[\widehat{\bm{U}}^{\mathsf T}\bm{S}\widehat{\bm{U}}\right]+\operatorname{tr}\left[\bm{H}\left(\bm{I}-\widehat{\bm{U}}^{\mathsf T}\widehat{\bm{U}}\right)\right],
  \label{align:ch16-习题16.2Lagrange}
\end{align}
其中 $\widehat{\bm{U}}$ 是一个 $D\times(D-M)$ 大小的矩阵，其列由 $\bm{u}_j$ 给出。对 $\widehat{\bm{U}}$ 求 $\tilde{J}$ 的最小值并证明解满足 $\bm{S}\widehat{\bm{U}}=\widehat{\bm{U}}\bm{H}$。显然，一个可能的解是，$\widehat{\bm{U}}$ 的列是 $\bm{S}$ 的特征向量。在这种情况下，$\bm{H}$ 是一个包含对应特征值的对角矩阵。假设 $\bm{H}$ 是一个对称矩阵，并且通过使用特征向量将其展开，证明 $\bm{S}\widehat{\bm{U}}=\widehat{\bm{U}}\bm{H}$ 的通用解（对于给出的 $\tilde{J}$ 值）与 $\widehat{\bm{U}}$ 的列是 $\bm{S}$ 的特征向量时的具体解相同。因为这些解都是等效的，所以选择特征向量解是很方便的。
\end{exercise}

\begin{exercise}\label{exercise:ch16-03}
（$\star$）假设\eqref{align:ch16-高维归一化}中的特征向量 $\bm{v}_i$ 具有单位长度，验证由该式定义的特征向量 $\bm{u}_i$ 已归一化至单位长度。
\end{exercise}

\begin{exercise}\label{exercise:ch16-04}
（$\star$）假设在概率 PCA 模型中将零均值、单位协方差的潜空间分布[\eqref{align:ch16-先验}]替换为一般形式的 Gauss 分布 $\mathcal{N}(\bm{z}|\bm{m},\bm{\Sigma})$。通过重新定义模型的参数，证明这种替换会使得任意有效的 $\bm{m}$ 和 $\bm{\Sigma}$ 都无法改变观测变量边缘分布 $p(\bm{x})$ 的形式。
\end{exercise}

\begin{exercise}\label{exercise:ch16-05}
（$\star\star$）假设 $\bm{x}$ 是一个 $D$ 维的随机变量，具有 Gauss 分布 $\mathcal{N}(\bm{x}|\bm{\mu},\bm{\Sigma})$，并考虑由 $\bm{y}=\bm{A}\bm{x}+\bm{b}$ 给出的另一个 $M$ 维随机变量，其中 $\bm{A}$ 是一个 $M\times D$ 大小的矩阵。证明 $\bm{y}$ 也具有 Gauss 分布，并找出其均值和协方差。讨论此 Gauss 分布在 $M<D$、$M=D$ 和 $M>D$ 情况下的形式。
\end{exercise}

\begin{exercise}\label{exercise:ch16-06}
（$\star\star$）利用一般分布的均值和协方差在仿射变换下的公式，推导概率 PCA 模型中边缘分布 $p(\bm{x})$ 的结果[\eqref{align:ch16-边缘分布}，协方差见\eqref{align:ch16-边缘协方差}]。
\end{exercise}

\begin{exercise}\label{exercise:ch16-07}
（$\star$）绘制\cref{subsection:ch16-生成式模型}描述的概率 PCA 模型的有向图，并将观测变量 $\bm{x}$ 的各个组成部分明确显示为单独的节点。验证概率 PCA 模型具有与\cref{subsection:ch11-朴素Bayes}讨论的朴素 Bayes 模型相同的独立结构。
\end{exercise}

\begin{exercise}\label{exercise:ch16-08}
（$\star\star$）利用线性-Gauss 模型的结果，证明概率 PCA 模型的后验分布 $p(\bm{z}|\bm{x})$ 可以由\eqref{align:ch16-后验}给出。
\end{exercise}

\begin{exercise}\label{exercise:ch16-09}
（$\star$）验证对概率 PCA 模型的对数似然函数[\eqref{align:ch16-对数似然}]关于参数 $\bm{\mu}$ 进行最大化会得到结果 $\bm{\mu}=\bar{\bm{x}}$，其中 $\bar{\bm{x}}$ 是数据向量的均值。
\end{exercise}

\begin{exercise}\label{exercise:ch16-10}
（$\star\star$）通过计算对数似然函数[\eqref{align:ch16-对数似然}]对参数 $\bm{\mu}$ 的二阶导数，证明驻点 $\bm{\mu}=\bar{\bm{x}}$ 代表唯一的最大值。
\end{exercise}

\begin{exercise}\label{exercise:ch16-11}
（$\star\star$）证明当 $\sigma^2\rightarrow0$ 时，概率 PCA 模型的后验均值相当于对主子空间的正交投影，这与传统 PCA 模型的结果一致。
\end{exercise}

\begin{exercise}\label{exercise:ch16-12}
（$\star\star$）对于 $\sigma^2>0$，证明概率 PCA 模型中的后验均值会相对于正交投影向原点移动。
\end{exercise}

\begin{exercise}\label{exercise:ch16-13}
（$\star\star$）证明在概率 PCA 下，根据传统 PCA 的最小二乘投影成本，数据点的最优重建由下式给出：
\begin{align}
  \tilde{\bm{x}}=\bm{W}_{\text{ML}}\left(\bm{W}_{\text{ML}}^{\mathsf T}\bm{W}_{\text{ML}}\right)^{-1}\bm{M}\,\mathbb{E}[\bm{z}|\bm{x}] .
  \label{align:ch16-习题16.13重建}
\end{align}
\end{exercise}

\begin{exercise}\label{exercise:ch16-14}
（$\star$）概率 PCA 模型中协方差矩阵的独立参数数量由\eqref{align:ch16-参数计数}给出。证明对于 $M=D-1$，独立参数的数量与一般 Gauss 分布中的相同；而对于 $M=0$，独立参数的数量与各向同性协方差 Gauss 分布相同。
\end{exercise}

\begin{exercise}\label{exercise:ch16-15}
（$\star$）推导出\cref{subsection:ch16-因子分析}描述的因子分析模型中独立参数数量的计算表达式。
\end{exercise}

\begin{exercise}\label{exercise:ch16-16}
（$\star\star$）证明\cref{subsection:ch16-因子分析}描述的因子分析模型不随潜空间坐标旋转而变化。
\end{exercise}

\begin{exercise}\label{exercise:ch16-17}
（$\star\star$）考虑一个线性-Gauss 潜变量模型，它具有潜空间分布 $p(\bm{z})=\mathcal{N}(\bm{z}|\bm{0},\bm{I})$ 和观测变量的条件分布 $p(\bm{x}|\bm{z})=\mathcal{N}(\bm{x}|\bm{W}\bm{z}+\bm{\mu},\bm{\Phi})$，其中 $\bm{\Phi}$ 是任意对称且正定的噪声协方差矩阵。对数据变量进行非奇异线性变换 $\bm{x}\rightarrow\bm{A}\bm{x}$，其中 $\bm{A}$ 是一个 $D\times D$ 矩阵。如果 $\bm{\mu}_{\text{ML}}$、$\bm{W}_{\text{ML}}$ 和 $\bm{\Phi}_{\text{ML}}$ 代表原始未变换数据所对应的最大似然解，证明 $\bm{A}\bm{\mu}_{\text{ML}}$、$\bm{A}\bm{W}_{\text{ML}}$ 和 $\bm{A}\bm{\Phi}_{\text{ML}}\bm{A}^{\mathsf T}$ 代表变换后数据所对应的最大似然解。最后证明在以下两种情况下模型的形式得以保留：(i) $\bm{A}$ 是对角矩阵且 $\bm{\Phi}$ 也是对角矩阵，这对应于因子分析。变换后的 $\bm{\Phi}$ 仍然是对角矩阵，因此因子分析在各成分矩阵的变换下具有协变性。(ii) $\bm{A}$ 是正交矩阵且 $\bm{\Phi}$ 与单位矩阵成比例，$\bm{\Phi}=\sigma^2\bm{I}$，这对应于概率 PCA。变换后的 $\bm{\Phi}$ 仍与单位矩阵成比例，因此概率 PCA 在数据空间的轴旋转下具有协变性，就像传统 PCA 一样。
\end{exercise}

\begin{exercise}\label{exercise:ch16-18}
（$\star$）验证连续潜变量模型对数似然函数可以写成两项之和[\eqref{align:ch16-ELBO分解}]，这两项分别由\eqref{align:ch16-ELBO定义}和\eqref{align:ch16-KL}定义。这可以通过使用概率的乘积法则来完成，公式如下：
\begin{align}
  p(\bm{x},\bm{z}|\bm{w})=p(\bm{z}|\bm{x},\bm{w})p(\bm{x}|\bm{w}),
  \label{align:ch16-习题16.18乘积}
\end{align}
然后将 $p(\bm{x},\bm{z}|\bm{w})$ 代入\eqref{align:ch16-ELBO定义}即可。
\end{exercise}

\begin{exercise}\label{exercise:ch16-19}
（$\star$）证明对于一组独立同分布数据，证据下界（ELBO）可以采用\eqref{align:ch16-ELBO-iid}的形式。
\end{exercise}

\begin{exercise}\label{exercise:ch16-20}
（$\star\star$）绘制一个有向图，以表示概率 PCA 模型的离散混合模型，其中的每个概率 PCA 模型都有自己的 $\bm{W}$、$\bm{\mu}$ 和 $\sigma^2$ 值。然后绘制一个修改后的有向图，这些参数值在该有向图所表示的混合模型的各部分之间是共享的。
\end{exercise}

\begin{exercise}\label{exercise:ch16-21}
（$\star\star$）通过最大化\eqref{align:ch16-完全数据期望}给出的完整数据对数似然函数的期望来推导概率 PCA 模型的 M 步骤方程[\eqref{align:ch16-M步W}]和[\eqref{align:ch16-M步sigma}]。
\end{exercise}

\begin{exercise}\label{exercise:ch16-22}
（$\star\star\star$）概率 PCA 的一个优点是，它可以应用于某些值缺失的数据集，前提是这些值是随机缺失的。在这种情况下，推导出最大化概率 PCA 模型似然函数的 EM 算法。注意，$\{\bm{z}_n\}$ 和向量 $\{\bm{x}_n\}$ 的缺失数据值现在都是潜变量。证明在所有数据值都观测到的特殊情况下，这可以简化为推导概率 PCA 的 EM 算法（见\cref{subsection:ch16-PCA的EM算法}）。
\end{exercise}

\begin{exercise}\label{exercise:ch16-23}
（$\star\star$）设 $\bm{W}$ 是一个 $D\times M$ 矩阵，其列定义了嵌入 $D$ 维数据空间的 $M$ 维线性子空间，并且 $\bm{\mu}$ 是一个 $D$ 维向量。给定一个数据集 $\{\bm{x}_n\}$，其中 $n=1,\cdots,N$，我们可以使用一组 $M$ 维向量 $\{\bm{z}_n\}$ 的线性映射来近似数据点，使得 $\bm{x}_n$ 能被 $\bm{W}\bm{z}_n+\bm{\mu}$ 近似。相关的平方和重建成本由下式给出：
\begin{align}
  J=\sum_{n=1}^{N}\left\|\bm{x}_n-\bm{\mu}-\bm{W}\bm{z}_n\right\|^2 .
  \label{align:ch16-习题16.23重建成本}
\end{align}
首先证明关于 $\bm{\mu}$ 最小化 $J$ 会得到一个相似的表达式，其中 $\bm{x}_n$ 和 $\bm{z}_n$ 分别替换为零均值变量 $\bm{x}_n-\bar{\bm{x}}$ 和 $\bm{z}_n-\bar{\bm{z}}$，其中 $\bar{\bm{x}}$ 和 $\bar{\bm{z}}$ 表示样本均值。然后证明在 $\bm{W}$ 保持不变的情况下关于 $\bm{z}_n$ 最小化 $J$ 会产生 PCA EM 算法的 E 步骤[\eqref{align:ch16-PCA-E步}]，并且在 $\{\bm{z}_n\}$ 保持不变的情况下关于 $\bm{W}$ 最小化 $J$ 会产生 PCA EM 算法的 M 步骤[\eqref{align:ch16-PCA-M步}]。
\end{exercise}

\begin{exercise}\label{exercise:ch16-24}
（$\star\star$）推导因子分析 EM 算法的 E 步骤方程[\eqref{align:ch16-FA-E步z}和\eqref{align:ch16-FA-E步zz}]。注意，从\cref{exercise:ch16-26}的结果可知，参数 $\bm{\mu}$ 可以用样本均值 $\bar{\bm{x}}$ 替换。
\end{exercise}

\begin{exercise}\label{exercise:ch16-25}
（$\star\star$）给出因子分析模型的完整数据对数似然函数的期望的表达式，从而推导出相应的 M 步骤方程[\eqref{align:ch16-FA-M步W}和\eqref{align:ch16-FA-M步Psi}]。
\end{exercise}

\begin{exercise}\label{exercise:ch16-26}
（$\star\star$）通过考虑二阶导数，证明\cref{subsection:ch16-因子分析}讨论的因子分析模型的对数似然函数相对于参数 $\bm{\mu}$ 的唯一驻点是由\eqref{align:ch16-均值}定义的样本均值给出的。此外，证明这个驻点是一个最大值。
\end{exercise}

\end{document}
